\chapter{绪论}

\section{研究背景与意义}

全球导航卫星系统（Global Navigation Satellite System, GNSS）已在城市规划、交通管理、灾害监测、精准农业等领域得到广泛应用。随着智慧城市和自动驾驶技术的发展，用户对定位服务的精度与可靠性提出了更高要求。然而，在城市峡谷、密集建筑群等复杂场景中，GNSS 信号传播路径易受高层建筑遮挡与反射影响，产生严重的多径效应和非视距（Non-Line-of-Sight, NLOS）传播问题。研究表明，NLOS 信号可导致伪距测量误差达数十米甚至上百米，使单点定位误差远超可接受范围，严重制约了高精度定位服务在城市环境中的可用性。

GNSS 接收机传统上采用最小二乘法或卡尔曼滤波进行位置解算，其基本假设是所有可见卫星信号均为视距（Line-of-Sight, LOS）传播。当该假设被 NLOS 信号破坏时，定位解算将产生系统性偏差。NLOS 信号识别面临以下主要挑战：首先，多径效应与 NLOS 传播的物理机制不同但在实际接收环境中往往同时存在，难以严格区分；其次，不同城市的建筑物密度、高度分布差异显著，导致 NLOS 信号的统计特性存在地域差异；此外，卫星相对于接收机的位置不断变化，同一地点的 NLOS 模式随时间变化；最后，车载导航等应用场景要求低延迟的定位解算，对模型复杂度提出严格约束。

准确识别并剔除或降权处理 NLOS 信号，成为提升 GNSS 定位性能的关键前提。通过有效的 NLOS 检测与抑制，可在不增加硬件成本的前提下显著改善定位精度与系统可用性，具有重要的理论意义与工程应用价值。如何在保持模型轻量级的同时提升跨场景泛化能力，是本研究的核心目标。

\section{国内外研究现状}

本节围绕 GNSS NLOS 信号识别这一核心任务，梳理了相关技术的发展脉络，对国内外代表性研究工作进行总结。

\subsection{阈值判别法}

阈值判别法基于载噪比（C/N0）、信号强度、伪距残差等单一或多个观测量的固定阈值进行 LOS/NLOS 分类。其优势在于实现简单、计算开销低，可直接嵌入接收机固件。然而，阈值选取高度依赖经验，难以适应不同城市环境、天气条件及接收机型号的变化，泛化能力有限。

\subsection{3D 模型辅助法}

3D 模型辅助法利用城市三维建筑模型与卫星星历，通过几何可见性分析预判信号传播路径是否被遮挡。该类方法物理意义明确、判别精度高，但依赖详细且实时更新的城市地理信息系统（GIS）数据，部署成本高昂，难以在中小城市或发展中国家推广。

\subsection{传统机器学习方法}

以支持向量机（SVM）、随机森林（Random Forest）、梯度提升树（GBDT）等为代表，可融合载噪比、卫星高度角、方位角、伪距残差等多维特征进行分类。相比阈值法，机器学习方法具备更强的非线性建模能力，但仍依赖人工特征工程，特征选取与组合策略对最终性能影响显著，且难以捕捉特征间的时空耦合关系。

\subsection{深度学习方法}

近年来，全连接神经网络（DNN）、长短期记忆网络（LSTM）、Transformer 等模型被引入 NLOS 识别任务。深度学习方法可自动从原始观测数据中提取判别性特征，避免繁琐的手工特征设计，在多个公开数据集上取得优异性能。然而，现有深度学习模型普遍存在可解释性弱、对输入维度固定化假设过强、难以同时建模卫星空间分布与信号时序演化等问题。

现有代表性方法在特定场景下取得了较好的识别效果，但仍存在以下不足：首先，多数方法仅利用单历元特征，忽略了卫星信号的时序演化信息；其次，跨场景泛化能力不足，在一个地点训练的模型直接部署到另一地点时性能显著下降；此外，模型复杂度与实时性之间的平衡尚未得到充分解决。

\section{研究内容与创新点}

针对现有方法的局限性，本文提出一种基于时空特征融合与稀疏表示的 GNSS NLOS 识别模型（Spatiotemporal Cross-Attention Model, STCA）。该模型借鉴 Zeng 等人（2024）提出的时空交叉注意力架构，结合 GNSS 信号传播的物理特性进行针对性改进。

本文设计了时空双通道特征提取架构。针对 GNSS 观测数据的固有特性，构建空间输入通道与时间输入通道：空间通道捕获同一时刻多颗可见卫星的分布特征（高度角、方位角、载噪比、伪距残差），时间通道追踪单颗卫星在连续多个历元的信号演化趋势。双通道设计使模型能够同时利用空间分集与时间相关性信息。

本文引入了交叉注意力机制实现时空特征融合。通过交叉注意力模块，使空间特征与时间特征在嵌入空间进行双向交互，自动学习不同卫星、不同时刻特征的重要性权重，避免简单拼接或加权平均导致的信息损失。

本文采用稀疏表示层增强特征判别能力。在分类器前引入稀疏性约束，迫使模型学习更具判别性的紧凑特征表示，提升对噪声与异常值的鲁棒性。

\section{论文章节安排}

本文一共包含五个章节，各章节的主要内容如下：

第一章为绪论。扼要介绍了 GNSS NLOS 信号识别的研究背景、核心挑战和研究价值，并系统综述了国内外相关研究进展。

第二章为相关研究。本章围绕 GNSS NLOS 信号检测任务，梳理了相关技术的发展脉络，包括基于几何约束的检测方法、基于传统机器学习的方法和基于深度学习的方法，并分析了各类方法的优势与不足。

第三章为 STCA 模型设计。本章首先给出引言，概述模型设计动机。然后详细介绍数据预处理流程和模型框架，包括空间编码器、时序编码器、交叉注意力融合模块、稀疏表示层和二分类器的设计原理与数学表述，并分析模型复杂度。

第四章为实验结果与分析。本章首先给出引言，概述实验目标。然后介绍实验设置与评价指标，展示外域场景下的实验结果，包括训练过程、混淆矩阵、ROC 曲线和 PR 曲线分析。最后通过时间窗口大小和模块消融实验验证各核心模块的有效性。

第五章为总结与展望。本章总结全文的研究工作与主要贡献，包括时空双通道架构设计、交叉注意力融合机制、稀疏表示层增强等创新点。同时客观分析当前方法的局限性，并从数据扩充、域泛化方法、时序模型改进、多星座融合、可解释性增强等方面展望未来的研究方向。

\section{本章小结}

本章阐述了 GNSS NLOS 识别问题的研究背景与实际意义，分析了 NLOS 检测面临的主要挑战，包括多径效应与 NLOS 的区分困难、城市环境复杂性、卫星运动的时变性以及实时性约束。系统综述了国内外相关研究进展，包括阈值判别法、3D 模型辅助法、传统机器学习方法和深度学习方法，并分析了各类方法的优缺点。针对现有方法未能同时利用时空特征、跨场景泛化能力不足等问题，本文提出了基于时空特征融合与稀疏表示的 STCA 模型。下一章将围绕 GNSS NLOS 信号检测任务，系统综述基于几何约束、传统机器学习和深度学习的检测方法研究进展。
